\chapter{Engine Implementation}

\section{Python Module}

The Python engine is implemented following the idiomatic approach for numerical computing in this language: vectorization via NumPy. Due to the Global Interpreter Lock (GIL) in standard CPython, explicit multithreading using the \texttt{threading} library does not yield true parallelism for CPU-bound tasks. Instead, our implementation allocates memory for all Monte Carlo paths simultaneously in a single NumPy array, delegating the heavy algebraic operations to underlying optimized libraries. 

While this vectorized approach achieves maximum possible performance in pure Python, it inherently incurs a high spatial complexity of \(\mathcal{O}(N)\). To mitigate this and leverage bare-metal performance, Python is repositioned as a high-level orchestrator. The Python frontend initializes the required \texttt{pybind11}-exposed classes (such as \texttt{MarketData} and \texttt{EuropeanOption}) and formats the inputs appropriately before invoking the backend. \texttt{pybind11} leverages the Python buffer protocol (via \texttt{pybind11::array\_t}) to expose the underlying C++ memory pointers directly to the Python environment. By avoiding deep copies of massive arrays, this \textit{zero-copy architecture} ensures that the intensive computational load is fully delegated to the C++ core with minimal data transfer latency, keeping the Python layer clean, readable, and focused exclusively on data preparation and execution flow control.

\section{C++ Module}

In contrast, the C++ module avoids large memory allocations and is designed around an iterative, multithreaded architecture. Since C++ does not have a GIL, it leverages multiple native OS threads to process the simulation paths concurrently. Each thread computes paths sequentially and maintains a local accumulator, resulting in a constant spatial complexity \(\mathcal{O}(1)\) per thread. This paradigm not only maximizes CPU utilization but also allows the engine to handle a massive number of paths without exhausting the system's RAM.

The core of the numerical integration resides in the hot loop of the \texttt{MonteCarloPricer::simulate} method. At the hardware level, this method is compiled leveraging aggressive optimization flags such as \texttt{-O3} and \texttt{-ffast-math}. The \texttt{-ffast-math} directive is particularly critical: it relaxes strict adherence to IEEE 754 standards, permitting the compiler to reassociate floating-point operations and aggressively vectorize transcendental functions, such as the exponential \(\exp(\cdot)\) present in the geometric Brownian motion equation. This optimization drastically reduces the CPU cycles required per path, under the strict mathematical assumption that neither \texttt{NaN} nor \texttt{Inf} values will be generated during the simulation domain.

Furthermore, when simulating deeply Out-of-The-Money (OTM) options where price trajectories can tend toward zero, standard IEEE 754 arithmetic handles underflow by generating subnormal (denormal) numbers. Processing these denormals requires microcode execution on the CPU, causing massive performance penalties. To guarantee peak throughput, our environment explicitly mitigates this by enabling the hardware-level SSE/AVX control flags \textit{Flush-to-Zero (FTZ)} and \textit{Denormals-Are-Zero (DAZ)}. This ensures that any subnormal operations are aggressively truncated to zero, preserving the cycle latency in the hot loop without compromising the financial accuracy of the simulation.

The price dynamic at each step is calculated directly within this optimized loop. Following the risk-neutral valuation, the asset price \(S_{T}\) and its antithetic counterpart are computed simultaneously:
\[
    S_{T}^{(+)} = S_0 \exp\left[ \left( r - \frac{1}{2}\sigma^2 \right)T + \sigma \sqrt{T} Z \right]
\]
\[
    S_{T}^{(-)} = S_0 \exp\left[ \left( r - \frac{1}{2}\sigma^2 \right)T - \sigma \sqrt{T} Z \right]
\]
where \(Z \sim \mathcal{N}(0,1)\) is drawn from a thread-local instance of \texttt{std::mt19937\_64}. By computing both paths intrinsically, the engine halves the calls to the pseudo-random number generator and exploits the negative covariance to reduce estimator variance, translating the theoretical principles into efficient machine code.

\section{Compilation Pipeline}

To construct the hybrid architecture, the compilation and linking process is managed via \texttt{CMake}. The \texttt{CMakeLists.txt} configuration file serves as the definitive blueprint for generating the shared library.

The pipeline integrates \texttt{pybind11} as a core dependency, directly wrapping the C++ objects (\texttt{MarketData}, \texttt{EuropeanOption}, \texttt{MonteCarloPricer}) and binding them to a native Python extension (e.g., a \texttt{.so} file on Unix-like systems or a \texttt{.pyd} on Windows). 

Crucially, to ensure the \texttt{-ffast-math} optimizations discussed in Section 4.2 are physically realized at the hardware level, the build system injects aggressive compiler flags directly into the target:

\begin{lstlisting}[language=cmake, caption={CMake compiler flags configuration}]
target_compile_options(MonteCarloEngine PRIVATE 
    -O3 
    -ffast-math 
    -march=native
)
\end{lstlisting}

The inclusion of the \texttt{-march=native} directive is imperative for modern quantitative engines. It instructs the compiler to bypass generic x86\_64 instruction sets and instead generate machine code tailored specifically to the host CPU's architecture. This explicitly enables advanced Single Instruction, Multiple Data (SIMD) registers (such as AVX2 or AVX-512). By doing so, the processor can vectorize the intensive arithmetic loops, computing multiple floating-point operations simultaneously per clock cycle and guaranteeing that the compiled extension delivers optimal, deterministic brute-force performance.
